\chapter{全文总结与展望}


\section{研究成果总结}

本文围绕“基于任务优先级的人形机器人腿臂协同控制”这一核心问题，针对人形机器人在复杂环境中执行全身动作时所面临的高维非线性、多关节耦合、多接触切换等挑战，系统开展了从运动学建模、学习控制到技能组合与实物验证的研究工作，构建了一套从运动先验获取、全身控制学习到多技能融合与真实系统验证的完整技术框架。本文的主要研究内容与创新点总结如下：

（1）提出了一种面向全身运动的改进重定向方法，实现了关节空间与任务空间的协同优化。针对人体与人形机器人之间在结构与自由度上的差异，构建了基于几何关系的关节角解析模型，并在此基础上引入构型保持机制与多末端任务空间约束，形成统一的优化框架。同时，通过非均匀缩放映射与带旋转约束的逆运动学精修过程，使机器人在保证末端位置跟踪精度的同时维持合理关节构型。实验结果表明，该方法在精度、自然性与计算效率之间取得了良好平衡，为后续控制方法提供了高质量的运动先验。

（2）提出了一种基于强化学习的全身控制框架，实现了从运动学一致性到动力学可执行性的跨越。在重定向结果的基础上，构建了以参考动作为先验的端到端强化学习控制方法，通过统一策略网络直接生成关节控制信号，使机器人能够在复杂接触与动力学约束下稳定执行全身动作。在方法设计上，通过构建包含动力学状态与参考信息的观测空间，并结合关键刚体跟踪、姿态稳定性及运动一致性等约束设计奖励函数，引导策略学习自然且稳定的运动行为。同时，引入参考状态初始化（RSI）、提前终止（ET）以及自适应难度采样等机制，有效提升了训练效率与稳定性。实物实验结果表明，该方法在真实机器人平台上具有良好的鲁棒性与可部署性。

（3）提出了一种基于技能嵌入的分层全身控制方法，实现了多技能的统一表示与灵活组合。针对单一策略难以应对复杂任务的问题，构建了“高层决策—低层跟踪”的分层控制框架。在底层，通过专家策略蒸馏获得通用全身动作跟踪控制器；在高层，通过自回归条件变分自编码器（MVAE）构建低维潜在动作空间，实现运动技能的紧凑表示与生成。通过在潜空间中进行策略决策，并结合底层控制器执行，实现了多种技能之间的平滑过渡与自然组合。同时，引入任意姿态跌倒恢复策略，显著提升系统在复杂环境中的鲁棒性。实验结果表明，该方法在技能组合能力与泛化性能方面具有明显优势。

（4）构建了面向真实系统的人形机器人实验平台，并系统验证了所提出方法的有效性。基于自研 X02 与 X03 人形机器人平台，建立了统一的硬件与软件系统架构，并针对绳驱耦合传动特性引入执行器感知的驱动约束机制，有效缩小了仿真与现实之间的差异。在此基础上，对动作重定向方法、强化学习控制方法以及技能嵌入控制框架进行了系统的实物验证。实验结果表明，本文所提出的方法能够在真实机器人平台上实现稳定、自然且具有结构性的全身运动，并在复杂任务中表现出良好的鲁棒性与适应能力。

综上所述，本文从运动学建模、学习控制到分层决策与实物验证，系统解决了人形机器人腿臂协同控制中的关键问题，实现了从单一动作跟踪到多技能融合控制的逐步提升，为人形机器人在复杂环境中的全身控制提供了一种具有系统性与可扩展性的技术框架。

\section{未来研究展望}

尽管本文在全身运动建模与控制方面取得了一定进展，但仍有若干关键问题有待进一步深入研究。围绕这些问题，未来可从以下几个方面展开：

（1）构建运动学与动力学一体化的全身控制框架。本文采用“运动学重定向 + 学习控制”的分阶段方法，在实际系统中取得了良好效果，但两者之间仍存在一定程度的解耦。在高动态动作或复杂接触任务中，运动学先验与动力学约束之间的潜在不一致性可能影响控制性能。未来可将运动学重定向与全身动力学控制方法（如层级QP或全身力控制）进行深度融合，实现从动作生成到力控制的统一建模，从而提升系统在复杂动态环境中的稳定性与可控性。

（2）提升控制策略的可解释性与任务优先级表达能力。强化学习方法在处理复杂动力学问题方面具有优势，但其内部决策机制缺乏明确的结构表达。在多任务与多技能场景下，策略如何隐式体现任务优先级与协同关系仍有待进一步分析。未来可结合任务优先级理论，在学习框架中引入显式结构约束或分层表示方法，从而增强策略的可解释性与稳定性。

（3）提升多技能统一控制与泛化能力。虽然本文通过技能嵌入方法实现了多种动作的组合与切换，但当前实验任务类型仍相对有限。在更加复杂或开放环境中，系统对未知任务的适应能力仍需进一步提升。未来可结合更大规模动作数据与生成模型，构建更加通用的技能表示空间，使机器人能够实现连续、多样且具有语义结构的行为生成。

（4）加强感知驱动的闭环控制能力。本文主要基于理想状态信息开展研究，而在真实系统中，感知噪声、时延及外部扰动等因素会对控制性能产生显著影响。未来可将视觉感知、力觉反馈等多模态信息引入控制框架，实现从“离线动作驱动”向“在线感知驱动”的转变，从而提升系统在动态环境中的实时决策与交互能力。

（5）推进仿真到现实的高效迁移方法研究。尽管本文通过执行器约束建模在一定程度上缩小了仿真与现实之间的差异，但在更复杂任务场景中仍可能存在性能退化问题。未来可结合系统辨识以及执行器精细建模等方法，进一步提升控制策略在真实系统中的泛化能力与部署可靠性。

总体而言，随着人形机器人技术的不断发展，如何实现高效、稳定且具有泛化能力的全身控制，将成为未来的重要研究方向。本文的研究工作为这一方向提供了有益探索，也为后续研究奠定了基础。